Figure 1 (frame size $100^\circ \times 70^\circ$) shows a part of the sky,
for 22 October 2009 at 21:00 local time. Four bright stars in the Perseus
and Andromeda constellations are missing in this chart. Find these missing
stars by looking at the sky. Then, draw a cross on the location of each
missing bright star in these two constellations on the chart (i.e.\
Figure 1). Use the numbers in Table 1-1 to indicate these crosses.

Note: Polaris is indicated by the ``N'' symbol in Figure 1.

\begin{center}
\begin{tabular}{ccl}
\multicolumn{3}{c}{Table 1-1}\\
Number & Common Name & Bayer Name \\
\hline
1 & Mirfak & Alpha Persei \\
2 & Alpheratz & Alpha Andromedae \\
3 & -- & Epsilon Persei \\
4 & Menkib & Xi Persei \\
5 & -- & Gamma Persei \\
6 & Algol & Beta Persei \\
7 & Almach & Gamma Andromedae \\
8 & -- & Delta Andromedae \\
9 & -- & 51 Andromedae \\
10 & Mirach & Beta Andromedae \\
11 & Atik & Zeta Persei \\
\end{tabular}
\end{center}

% FIGURE NEEDED: Figure 1, star chart (100 x 70 degrees) of the
% Perseus--Andromeda region with four bright stars missing, Polaris marked N
